\pdfvariable trailerid {[<C2992026090200000000000000000000><C2992026090200000000000000000000>]}
\documentclass[10pt]{article}
\usepackage[margin=0.76in]{geometry}
\usepackage{amsmath,amssymb,amsthm,booktabs,microtype,hyperref}
\hypersetup{colorlinks=true,linkcolor=blue,urlcolor=blue}
\urlstyle{same}
\providecommand{\CRevisionRound}{2}
\newtheorem{theorem}{Theorem}
\newtheorem{proposition}{Proposition}
\newcommand{\R}{\mathbb R}
\newcommand{\dd}{\,\mathrm d}
\newcommand{\Ei}{\operatorname{Ei}}
\setlength{\emergencystretch}{3em}
\title{The Radial Lamb--Oseen Similarity Class:\\
Uniqueness, Exact Particle Angles, and Dissipation Boundaries}
\author{Route-A source-local certificate HCS-C299}
\date{2 September 2026}

\begin{document}
\maketitle

\begin{abstract}
We classify the bounded-at-origin, finite-circulation, radial forward-self-
similar solutions of the planar viscous vorticity equation: the signed
Lamb--Oseen Gaussian is the entire declared class.
\ifnum\CRevisionRound=0
We reconstruct its velocity, integrate every positive-radius particle path,
and give all radial moments, finite $L^p$ norms, enstrophy, and palinstrophy.
\fi
\ifnum\CRevisionRound>0
Beyond exact Eulerian and Lagrangian formulas, we separate the zero-age,
inviscid, origin-particle, long-time, recurrence, and infinite-energy
boundaries without extending uniqueness to arbitrary vortices.
\fi
\ifnum\CRevisionRound>1
A 213-cell evidence certificate and independent symbolic, replay, mutation,
and deterministic-build lanes support regression; the analytic proof remains
global.  The resulting Route-A tuple is five explicit failures.
\fi
\end{abstract}

\section{Declared class and complete classification}
On $\R^2$ consider
\begin{equation}\label{eq:NS}
 \partial_t\omega+u\cdot\nabla\omega=\nu\Delta\omega,
 \qquad \nabla\cdot u=0,\qquad \operatorname{curl}u=\omega,
\end{equation}
with the decaying planar Biot--Savart velocity.  Fix
$\Gamma\in\R$, $\nu>0$, and $\tau=t+\tau_0>0$, where $\tau_0\ge0$.

\begin{theorem}[Radial forward-similar uniqueness]\label{thm:class}
Suppose a classical radial solution of~\eqref{eq:NS} has
\begin{equation}\label{eq:ansatz}
 \omega(x,t)=\tau^{-1}F(\xi),\qquad \xi=|x|/\sqrt\tau,
\end{equation}
where $F\in C^2([0,\infty))$ is bounded at the origin and
$\int_0^\infty |F(\xi)|\xi\dd\xi<\infty$.  If
$\int_{\R^2}\omega\dd x=\Gamma$, then
\begin{equation}\label{eq:omega}
 \omega(x,t)=\frac{\Gamma}{4\pi\nu\tau}
 \exp\!\left(-\frac{|x|^2}{4\nu\tau}\right).
\end{equation}
Thus~\eqref{eq:omega} is the entire declared similarity class, including
$\Gamma=0$.
\end{theorem}

\begin{proof}
Radial vorticity has radial gradient, whereas its Biot--Savart velocity is
tangential; hence $u\cdot\nabla\omega=0$.  Substitution of~\eqref{eq:ansatz}
into the heat equation yields
\begin{equation}\label{eq:ode}
 \nu(F''+\xi^{-1}F')+F+\frac{\xi}{2}F'=0.
\end{equation}
Multiplying by $\xi$ and integrating once gives
\begin{equation}\label{eq:first}
 \nu\xi F'(\xi)+\frac{\xi^2}{2}F(\xi)=C.
\end{equation}
The origin hypotheses make both left-hand terms tend to zero, so $C=0$.
Therefore $F=Ae^{-\xi^2/(4\nu)}$.  The normalization
$\Gamma=2\pi\int_0^\infty F(\xi)\xi\dd\xi=4\pi\nu A$ fixes $A$ and proves
\eqref{eq:omega}.
\end{proof}

The scope of Theorem~\ref{thm:class} is essential: it classifies every
profile satisfying~\eqref{eq:ansatz} and the displayed regularity and
integrability conditions, not arbitrary vortex filaments or arbitrary
solutions of~\eqref{eq:NS}.  The classical Lamb--Oseen mechanism is source
owned by Oseen; we claim no historical priority.

\section{Velocity and exact Lagrangian dynamics}
Polar circulation reconstructs the velocity without an unrecorded constant:
\begin{equation}\label{eq:velocity}
 u_r=0,\qquad
 u_\theta(r,t)=\frac{\Gamma}{2\pi r}
 \left(1-e^{-r^2/(4\nu\tau)}\right)\quad(r>0),
 \qquad u(0,t)=0.
\end{equation}
Indeed $2\pi r u_\theta=2\pi\int_0^r s\omega(s,t)\dd s$; near zero,
$u_\theta=\Gamma r/(8\pi\nu\tau)+O(r^3)$.

\begin{proposition}[Every particle orbit]\label{prop:particle}
For $r_0>0$, radius is invariant.  Put $a=r_0^2/(4\nu)$ and
\begin{equation}\label{eq:primitive}
 \mathcal F_a(\tau)=\tau-\tau e^{-a/\tau}-a\Ei(-a/\tau).
\end{equation}
Between positive ages $\tau_s<\tau_t$,
\begin{equation}\label{eq:angle}
 \theta(\tau_t)-\theta(\tau_s)
 =\frac{\Gamma}{2\pi r_0^2}
 [\mathcal F_a(\tau_t)-\mathcal F_a(\tau_s)].
\end{equation}
Here $\theta$ denotes a continuous real-valued lift of polar angle, not only
its class modulo $2\pi$.
The origin is a fixed particle.  If $\Gamma\ne0$ and $r_0>0$, then
\begin{equation}\label{eq:angleasymp}
 \theta(t)=\frac{\Gamma}{8\pi\nu}\log\tau+O(1)
 \qquad(\tau\to\infty).
\end{equation}
\end{proposition}
\begin{proof}
Equation~\eqref{eq:velocity} gives $\dot r=0$ and
\[
 \dot\theta=\frac{\Gamma}{2\pi r_0^2}(1-e^{-a/\tau}).
\]
Direct differentiation shows $\mathcal F_a'=1-e^{-a/\tau}$, proving
\eqref{eq:angle}.  The continuous origin value proves the fixed-point claim.
Finally $1-e^{-a/\tau}=a/\tau+O(\tau^{-2})$, which gives
\eqref{eq:angleasymp} because $a/r_0^2=1/(4\nu)$.
\end{proof}

\section{All moments, norms, and exact dissipation}
The change of variable $y=r^2/(4\nu\tau)$ evaluates the whole parameter
family, rather than a finite selection.
\begin{proposition}\label{prop:norms}
For every integer $k\ge0$ and every real $1\le p<\infty$,
\begin{align}
 \int_{\R^2}|x|^{2k}\omega(x,t)\dd x
 &=\Gamma k!(4\nu\tau)^k,\label{eq:moments}\\
 \|\omega(\cdot,t)\|_p^p
 &=\frac{|\Gamma|^p}{p(4\pi\nu\tau)^{p-1}}.\label{eq:lp}
\end{align}
Also $\|\omega\|_\infty=|\Gamma|/(4\pi\nu\tau)$ and
\begin{align}
 \int_{\R^2}\omega^2\dd x&=\frac{\Gamma^2}{8\pi\nu\tau},
 \label{eq:enstrophy}\\
 \int_{\R^2}|\nabla\omega|^2\dd x&=
 \frac{\Gamma^2}{16\pi\nu^2\tau^2},\label{eq:palinstrophy}\\
 \frac{\dd}{\dd t}\int_{\R^2}\omega^2\dd x
 &=-2\nu\int_{\R^2}|\nabla\omega|^2\dd x.\label{eq:dissipation}
\end{align}
\end{proposition}
\begin{proof}
Polar integration and $y=r^2/(4\nu\tau)$ reduce~\eqref{eq:moments} to
$\int_0^\infty y^k e^{-y}\dd y=k!$ and~\eqref{eq:lp} to
$\int_0^\infty e^{-py}\dd y=1/p$.  The supremum occurs at the origin.
The case $p=2$ gives~\eqref{eq:enstrophy}; differentiating~\eqref{eq:omega}
radially and integrating gives~\eqref{eq:palinstrophy}.  Differentiation in
$t$ closes~\eqref{eq:dissipation}.
\end{proof}

\ifnum\CRevisionRound>0
\section{Boundary atlas and recurrence obstruction}
The formulas remain useful only if singular endpoints are not silently
substituted into a positive-parameter theorem.
\begin{center}
\small
\begin{tabular}{@{}llp{0.62\linewidth}@{}}
\toprule
Boundary & Status & Exact conclusion \\
\midrule
$\Gamma=0$ & included & $\omega=u=0$ at every positive age. \\
$\tau_0>0$ & included & Smooth finite-enstrophy initial profile. \\
$\tau_0=0$ & weak trace & $\omega(\cdot,t)\rightharpoonup\Gamma\delta_0$ as $t\downarrow0$; smooth for $t>0$. \\
$\nu\downarrow0$ & weak limit & Vorticity tends to $\Gamma\delta_0$; velocity tends off the origin to the point vortex. \\
$r_0=0$ & separate & The continuous velocity fixes the origin; formula~\eqref{eq:angle} is only for $r_0>0$. \\
$\tau\to\infty$ & asymptotic & Positive-radius angular drift is logarithmic as in~\eqref{eq:angleasymp}. \\
$p>1$ & obstruction & For $\Gamma\ne0$, every finite $L^p$ norm decreases strictly with age. \\
$R\to\infty$ & warning & Whole-plane kinetic energy diverges logarithmically. \\
\bottomrule
\end{tabular}
\end{center}
Positive viscosity is the theorem domain; $\nu=0$ is never inserted into
\eqref{eq:omega}.  The $L^p$ law implies that no nonzero vorticity state in
this family can recur at two distinct ages.  This fluid-state statement does
not mislabel the nonautonomous angular motion as a primitive periodic orbit.

The far field $u_\theta=\Gamma/(2\pi r)+o(r^{-1})$ yields the sharp disk
asymptotic
\begin{equation}\label{eq:energy}
 \frac12\int_{|x|<R}|u(x,t)|^2\dd x
 =\frac{\Gamma^2}{4\pi}\log R+O(1).
\end{equation}
Thus positive-age enstrophy and palinstrophy are finite, whereas nonzero
whole-plane kinetic energy is not.  These are different integrability tests.
\fi

\ifnum\CRevisionRound>1
\section{Exact evidence and hostile audit}
The archived evidence covers eight signed field cases.  Each contains nine
similarity coordinates, nine even-moment receipts, and six $L^p$ receipts.
Twelve positive-radius cases compare~\eqref{eq:primitive} against independent
90-digit quadrature, and nine rows encode the boundary ledger: in total
$72+72+48+12+9=213$ audited cells.  This is regression evidence only; the
global statements follow from the proofs above.

The independent checker imports no producer and reconstructs every rational
or transcendental value.  A separate SymPy lane verifies~\eqref{eq:ode},
the heat equation, curl, all sampled moments and norms, the dissipation
identity, the exponential-integral derivative, and both asymptotic
coefficients.  Byte replay, repaired-hash hostile mutations, duplicate-key
and nonfinite-token rejection, exact YAML semantics, and six isolated
fixed-epoch builds close the digital audit.

\paragraph{Collision boundaries.}
C206 concerns Couette advection--diffusion and Fourier shearing on
$\mathbb T\times\R$; here the planar velocity is reconstructed by
Biot--Savart and radial geometry cancels nonlinear advection.  C207 concerns
scalar nonlinear diffusion and Barenblatt profiles; here diffusion is linear
only after the vorticity reduction, and circulation, particle trajectories,
the point-vortex boundary, and kinetic energy are part of the theorem.

\section{Route-A verdict and reproducibility}
Under the frozen scope \texttt{NO\_BAD\_EULER\_OR\_ROOT\_NUMBER}, the exact
tuple is
\[
 (\mathrm{A0\_FAIL},\mathrm{A1\_FAIL},\mathrm{A2\_FAIL},
   \mathrm{A3\_FAIL},\mathrm{A4\_FAIL}).
\]
Circulation and Gaussian moments are not arithmetic local data (A0).  Strict
$L^p$ decay supplies no recurrent primitive-orbit bridge (A1).  Physical
viscous time, even with logarithmic angular drift, is not an arithmetic clock
(A2).  No target completed function or functional equation appears (A3).
The dissipative heat/Navier--Stokes generator is not a self-adjoint
Hilbert--P\'olya operator (A4).  The overall verdict is
\texttt{ROUTE\_A\_REJECTED}, and Route B stays locked.

No target Euler factor, root number, automorphy statement, divisor law,
functional equation, target-zero correspondence, or target spectral
realization is claimed.

\paragraph{Reproducibility statement.}
The release retains the producer, duplicate-rejecting checker, symbolic lane,
byte replay, mutation suite, strict evaluation, three manuscript revisions,
and self-excluding manifest.  Source commit, evaluator digest, epoch, and
scope are machine checked.

\paragraph{AI-use statement.}
A generative language model assisted with drafting and verification-code
scaffolding.  Independent scripted lanes checked formulas, exact outputs,
scope boundaries, and release artifacts; final responsibility remains with
the authors.
\fi

\section*{Source lineage}
Oseen's 1912 line-vortex work is the classical owner token for the profile.
Gallay and Wayne provide modern mathematical context for the two-dimensional
Oseen vortex and its stability.  These citations establish lineage, not a
novelty or priority claim for the classical formula.

\begin{thebibliography}{9}
\bibitem{Oseen1912}
C. W. Oseen, ``Über Wirbelbewegung in einer reibenden Flüssigkeit,''
\emph{Arkiv för Matematik, Astronomi och Fysik} 7 (1912), no.~14.

\bibitem{GallayWayne2005}
T. Gallay and C. E. Wayne, ``Global stability of vortex solutions of the
two-dimensional Navier--Stokes equation,''
\emph{Communications in Mathematical Physics} 255 (2005), 97--129.
DOI: \href{https://doi.org/10.1007/s00220-004-1254-9}{10.1007/s00220-004-1254-9}.
\end{thebibliography}

\end{document}
